\subsection{分裂半单李代数}

定义 1. 设 \(\mathfrak{g}\) 是半单李代数。嘉当子代数 \(\mathfrak{h}\) （\(\mathfrak{g}\) 的）称为分裂的，若对一切 \(x \in \mathfrak{h}\) ，\(\operatorname{ad}_{\mathfrak{g}} x\) 可三角化。半单李代数称为可分裂的，若它有一个分裂嘉当子代数。分裂半单李代数是指一个偶 \((\mathfrak{g}, \mathfrak{h})\) ，其中 \(\mathfrak{g}\) 是半单李代数，而 \(\mathfrak{h}\) 是 \(\mathfrak{g}\) 的分裂嘉当子代数。

注. 1) 设 \(\mathfrak{g}\) 是半单李代数，\(\mathfrak{h}\) 是 \(\mathfrak{g}\) 的嘉当子代数。对一切 \(x \in \mathfrak{h}\) ，\(\operatorname{ad}_{\mathfrak{g}} x\) 半单（第七章 §2 第 4 节 定理 2）。于是，说 \(\mathfrak{h}\) 分裂就意味着 \(\operatorname{ad}_{\mathfrak{g}} x\) 对一切 \(x \in \mathfrak{h}\) 可对角化。

\begin{enumerate}
\def\labelenumi{\arabic{enumi})}
\setcounter{enumi}{1}
\item
  若 \(k\) 代数闭，则每个半单李代数 \(\mathfrak{g}\) 都可分裂，且 \(\mathfrak{g}\) 的每个嘉当子代数都分裂。当 \(k\) 不是代数闭时，存在不可分裂的半单李代数（习题 2 a)）；此外，若 \(\mathfrak{g}\) 可分裂，\(\mathfrak{g}\) 也可能有不分裂的嘉当子代数（习题 2 b)）。
\item
  设 \(\mathfrak{g}\) 是半单李代数，\(\mathfrak{h}\) 是 \(\mathfrak{g}\) 的嘉当子代数，\(\rho\) 是 \(\mathfrak{g}\) 的有限维单射表示且使 \(\rho(\mathfrak{h})\) 可对角化。则 \(\operatorname{ad}_{\mathfrak{g}}x\) 对一切 \(x \in \mathfrak{h}\) 可对角化（第七章 §2 第 1 节 例 2），故 \(\mathfrak{h}\) 分裂。
\item
  我们将看到（§3 第 3 节 命题 10 的推论），若 \(\mathfrak{h}\) 与 \(\mathfrak{h}'\) 是 \(\mathfrak{g}\) 的分裂嘉当子代数，则存在 \(\mathfrak{g}\) 的初等自同构把 \(\mathfrak{h}\) 变为 \(\mathfrak{h}'\) 。
\item
  设 \(\mathfrak{g}\) 是约化李代数。则 \(\mathfrak{g} = \mathfrak{c} \times \mathfrak{s}\) ，其中 \(\mathfrak{c}\) 是 \(\mathfrak{g}\) 的中心，而 \(\mathfrak{s} = \mathcal{D}\mathfrak{g}\) 半单。\(\mathfrak{g}\) 的嘉当子代数正是形如 \(\mathfrak{h} = \mathfrak{c} \times \mathfrak{h}'\) 的子代数，其中 \(\mathfrak{h}'\) 是 \(\mathfrak{s}\) 的嘉当子代数（第七章 §2 第 1 节 命题 2）。此时称 \(\mathfrak{h}\) 分裂，若 \(\mathfrak{h}'\) 相对于 \(\mathfrak{s}\) 分裂。这以一种明显的方式引出可分裂或分裂约化代数的定义。
\end{enumerate}
% semantic-fragments: legacy-input-seam
\relax{}
